
\keywords{linear logic, parallel programming, logic programming}

\begin{center}
{\huge \bf 
   Linear Logic and Coordination for Parallel Programming
}
\vspace{.25in}
{\LARGE
\\Fl\'{a}vio Manuel Fernandes Cruz\\}
\vspace{.25in}
{\large 
  CMU-CS-15-148
\\ March 2016 \\}
 \vspace{.5in}

\setlength{\tabcolsep}{2em}
{\large
\begin{tabular}{cc}
\\ Computer Science Department & Consortium of:   \\
School of Computer Science & Universidades do Minho,  Aveiro e Porto \\
Carnegie Mellon University & Portugal \\
Pittsburgh PA, USA &  \\
\end{tabular}
}

\begin{tabular}{cccc}
\\ \includegraphics[width=1.2in]{logos/cmu.eps} & \includegraphics[width=1.0in]{logos/um.jpg}  & \includegraphics[width=1.0in]{logos/ua.jpg} & \includegraphics[width=1.2in]{logos/up.eps} \\
\end{tabular}

%\vspace{0.3in}

{\bf Thesis Committee \\}

Frank Pfenning, Carnegie Mellon University (Co-Chair)\\
Seth Goldstein, Carnegie Mellon University (Co-Chair)\\
Ricardo Rocha, University of Porto (Co-Chair)\\
Umut Acar, Carnegie Mellon University\\
Luis Barbosa, University of Minho\\
Carlos Guestrin, University of Washington\\\
\vspace{0.2in}

\emph{Submitted in partial fulfillment of the requirements for the degree of Doctor of Philosophy}

\vspace{0.1in}
%\vfill

Copyright \copyright\ 2016 Fl\'{a}vio Cruz

\vspace{0.2in}

{\medskip\footnotesize\noindent
Support for this research was provided by the European Regional Development
Fund (ERDF) through the Operational Programme for Competitiveness and
Internationalisation - COMPETE 2020 Programme; by the Portuguese funding agency
- Funda\c{c}\~{a}o para a Ci\^{e}ncia e a
Tecnologia (Portuguese Foundation for Science and Technology) through the
Carnegie Mellon Portugal Program, under grants SFRH / BD / 51566 / 2011 and
within projects POCI-01-0145-FEDER-006961 and FCOMP-01-0124-FEDER-037281.

The views and conclusions contained in this document are those of the author and
should not be interpreted as representing the official policies, either
expressed or implied, of any sponsoring institution, the U.S. government or any
other entity.
\par}

\end{center}
